\documentclass[conference]{IEEEtran}
\IEEEoverridecommandlockouts
\usepackage{cite}
\usepackage{amsmath,amssymb,amsfonts}
\usepackage{algorithmic}
\usepackage{graphicx}
\usepackage{textcomp}
\usepackage{xcolor}
\usepackage[utf8]{inputenc}
\usepackage{hyperref}

\def\BibTeX{{\rm B\kern-.05em{\sc i\kern-.025em b}\kern-.08em
    T\kern-.1667em\lower.7ex\hbox{E}\kern-.125emX}}
\begin{document}

\title{LocalDemand: Final Literature Review}

\author{\IEEEauthorblockN{B.Tech Capstone Project}
\IEEEauthorblockA{\textit{LocalDemand}\\
}
}

\maketitle

\begin{abstract}
This literature review synthesizes 25 research papers relevant to demand forecasting, small and medium enterprise (SME) technology adoption, and inventory management. The papers are categorized into six thematic sections that address the core challenges of our proposed application, LocalDemand. Each section highlights how the reviewed literature supports our approach while identifying the gaps that LocalDemand aims to fill.
\end{abstract}

\begin{IEEEkeywords}
Demand Forecasting, SME, Inventory Management, LocalDemand
\end{IEEEkeywords}

\section{Forecasting Methods}

\subsection{Comparative Evaluation of SARIMA, FB Prophet, and LSTM for Hourly Demand Forecasting in Saudagar Kopi}
\textbf{Author(s) \& Year:} Kevina Syahla Aqila et al. (2026)\\
\textbf{Model/Method used:} SARIMA, FB Prophet, Stacked LSTM.\\
\textbf{Dataset used:} Hourly POS sales data from Saudagar Kopi in Jakarta.\\
\textbf{Results:} LSTM RMSE = 0.1535, FB Prophet RMSE = 0.46, SARIMA RMSE = 0.5265.\\
\textbf{Connection to our PS / gap:} Strongly relates to the forecast-to-action gap (3). It successfully predicts volume but does not convert insights into automated inventory decisions, which LocalDemand addresses.\\
\textbf{Literature Review Paragraph:}\\
Recent advancements in demand forecasting have increasingly targeted the operational inefficiencies of small food and beverage businesses. Aqila et al. (2026) conducted a comparative empirical evaluation of SARIMA, FB Prophet, and LSTM models to predict high-frequency, hourly demand at a Jakarta coffee shop. The authors incorporated exogenous variables, such as rush-hour and holiday indicators, to handle significant demand volatility. Their results revealed that while LSTM provided the highest accuracy (RMSE = 0.1535), FB Prophet also produced highly stable and interpretable forecasts (RMSE = 0.46) by effectively managing weekly and daily seasonality. Despite demonstrating that complex time-series models can successfully anticipate peak congestion in SME environments, the study remains purely analytical. It successfully predicts customer volume but fails to translate these insights into automated inventory decisions. LocalDemand addresses this exact limitation by utilizing a Prophet-based architecture not just to forecast sales, but to explicitly generate actionable daily restock recommendations.

\subsection{Smart Sales Forecasting Machine Learning Models for Demand Prediction in Retail}
\textbf{Author(s) \& Year:} Sandeep V et al. (2025)\\
\textbf{Model/Method used:} Hierarchical Agglomerative Clustering (HAC), K-Means, Decision Trees, Linear Regression.\\
\textbf{Connection to our PS / gap:} Validates the integration of external regressors (weather) to improve accuracy.\\
\textbf{Literature Review Paragraph:}\\
Traditional forecasting methods that rely solely on historical sales data frequently fail to capture the nuances of consumer purchasing behavior, which is heavily influenced by the external environment. Sandeep et al. (2025) addressed this by developing a machine learning framework that incorporates external variables---such as weather conditions and economic indicators---into retail demand prediction. Utilizing K-Means and Hierarchical Agglomerative Clustering, the authors grouped products with similar demand profiles before applying decision trees and linear regression. Their findings indicate that augmenting historical data with environmental regressors significantly enhances forecasting adaptability and accuracy. While the study provides a strong methodological justification for using external data, its reliance on established, static datasets (e.g., Big Mart Sales) limits its applicability to the highly volatile, data-scarce reality of micro-SMEs. LocalDemand adopts the core philosophy of this research---specifically the integration of weather and holiday variables---but executes it using Facebook Prophet to explicitly handle the data scarcity and cold-start problems that standard regression and clustering algorithms struggle to manage.

\subsection{State-Level Drug Retail Sales Forecasting Using XGBoost and Facebook Prophet}
\textbf{Author(s) \& Year:} Mohit Kamboj et al. (2025)\\
\textbf{Model/Method used:} Hybrid Facebook Prophet + XGBoost.\\
\textbf{Results:} Reduced median state-wise MAPE by 12\% to 18\%.\\
\textbf{Connection to our PS / gap:} Empirically validates Prophet's superiority in retail context handling seasonality and holidays.\\
\textbf{Literature Review Paragraph:}\\
The increasing complexity of modern retail environments has led to the adoption of hybrid machine learning architectures that combine statistical interpretability with nonlinear pattern recognition. Kamboj et al. (2025) demonstrated this by developing a hybrid forecasting framework using Facebook Prophet and XGBoost to predict state-level sales for the Rossmann drugstore chain in Germany. Prophet was utilized to achieve robust time-series decomposition (handling holidays and seasonality), while XGBoost was applied to the residuals to capture the nonlinear effects of internal promotions and competitor proximity. The hybrid model outperformed classical baselines like SARIMA and Holt-Winters, reducing the median state-wise MAPE by 12\% to 18\%. This study provides powerful empirical validation for Facebook Prophet's efficacy in retail demand forecasting, particularly regarding its handling of exogenous events like holidays. However, the study's focus on macro-level (state-wide) forecasting for a massive corporate enterprise means its findings are not directly actionable for single-location SMEs. LocalDemand scales this advanced Prophet-based methodology down to the micro-level, applying it directly to individual menu items to generate accessible, localized restock recommendations for small food businesses.

\subsection{Automated Demand Forecasting in Small to Medium-Sized Enterprises}
\textbf{Author(s) \& Year:} Thomas Gärtner et al. (2024)\\
\textbf{Model/Method used:} ARIMA, SARIMAX, Holt-Winters, XGBoost, Dilated CNN, GAM.\\
\textbf{Results:} Dilated CNN achieved nRMSE = 0.3119, MAPE = 33.19\%.\\
\textbf{Literature Review Paragraph:}\\
Gärtner et al. (2024) presented an automated demand forecasting system to help small and medium-sized enterprises (SMEs) improve their inventory planning and reduce manual forecasting efforts. The study used several forecasting models, including ARIMA, SARIMAX, Holt-Winters, XGBoost, Dilated CNN, and Ensemble models, to identify the most suitable forecasting approach. The results showed that the Dilated CNN model produced the most accurate forecasts and performed better than traditional forecasting methods in most of the SMEs included in the study. The proposed system helps businesses make better inventory decisions without requiring advanced technical knowledge. However, the model mainly depends on historical sales data and does not consider external factors such as weather or holidays. It also requires sufficient historical data, making it less suitable for newly established businesses.

\subsection{Enhancing Retail Sales Forecasting with Optimized Machine Learning Models}
\textbf{Author(s) \& Year:} Priyam Ganguly, Isha Mukherjee (2024)\\
\textbf{Model/Method used:} Random Forest, Linear Regression, Gradient Boosting, SVR, XGBoost.\\
\textbf{Results:} Optimized Random Forest achieved R$^2$ score of 0.945.\\
\textbf{Literature Review Paragraph:}\\
Ganguly and Mukherjee (2024) proposed a machine learning-based approach to improve retail sales forecasting and inventory planning. The study compared different machine learning models such as Linear Regression, Random Forest, Gradient Boosting, Support Vector Regression, and XGBoost using a retail sales dataset. Hyperparameter tuning was applied to improve the performance of the models. The experimental results showed that the optimized Random Forest model achieved the highest forecasting accuracy compared to the other models. The study concluded that machine learning techniques can effectively predict future sales and support better inventory management. However, the research mainly focused on retail stores and did not provide inventory recommendations or consider external factors like weather and holidays. It also did not develop a mobile application for small food businesses.

\subsection{Optimization of Forecasting Performance in the Retail Sector Using Artificial Intelligence}
\textbf{Author(s) \& Year:} Hoda Jatte et al. (2025)\\
\textbf{Model/Method used:} Linear Regression, Decision Tree, Random Forest, XGBoost, Prophet, LSTM.\\
\textbf{Results:} LSTM achieved 92.31\% accuracy, Prophet achieved 85.71\% accuracy.\\
\textbf{Literature Review Paragraph:}\\
Jatte et al. (2025) presented a comparative study on optimizing retail demand forecasting using Artificial Intelligence techniques to improve inventory management and supply chain operations. The authors employed six forecasting models, namely Linear Regression, Decision Tree, Random Forest, XGBoost, Prophet, and Long Short-Term Memory (LSTM), using a Kaggle retail demand forecasting dataset containing over 100,000 historical sales records. The methodology included data preprocessing, feature engineering, model training, and performance evaluation using Accuracy, Precision, Recall, F1-score, RMSE, and R$^2$ metrics. The experimental results demonstrated that the LSTM model achieved the highest forecasting performance with 92.31\% accuracy and 96\% F1-score, while the Prophet model also produced reliable forecasting results with 85.71\% accuracy. The study concluded that AI-based forecasting models significantly improve demand prediction and inventory planning in the retail sector. However, the research was limited to a public retail dataset and did not validate the proposed models on small food businesses such as bakeries, cafés, and restaurants. Furthermore, the study did not incorporate comprehensive external factors such as weather and holidays or provide a real-time mobile inventory recommendation system, leaving scope for future research.

\section{Small/Sparse Data Forecasting}

\subsection{Performance Analysis of Time Series Models to Predict Sales Activity of SMEs}
\textbf{Author(s) \& Year:} Justine Christian Susanto et al. (2026)\\
\textbf{Model/Method used:} XGBoost regression, 7-day Moving Average (MA(7)).\\
\textbf{Connection to our PS / gap:} Addresses data scarcity gap (1) focusing on zero-inflation which is common in small businesses.\\
\textbf{Literature Review Paragraph:}\\
Forecasting for small enterprises is frequently complicated by highly irregular demand patterns and data sparsity. Susanto et al. (2026) investigated the efficacy of time-series models in predicting sales for an SME operating under a pre-order framework, a system characterized by sporadic demand spikes and high frequencies of zero-transaction days (zero inflation). The authors benchmarked an XGBoost regression model against a simple 7-day Moving Average. They discovered that despite its sophistication, XGBoost struggled to accurately capture the discontinuous nature of the pre-order data, with the basic moving average providing more stable baseline estimates. The study concluded that continuous regression models are fundamentally challenged by zero-inflated, low-volume data environments. This directly highlights the exact data scarcity gap that LocalDemand addresses. By acknowledging that standard regression models fail under sparse conditions, LocalDemand utilizes Facebook Prophet---which inherently handles missing data points and trend shifts more gracefully---and incorporates external regressors to stabilize predictions even when historical daily sales volumes are intermittent or zero-inflated.

\subsection{LSTM-Based Hourly Sales Forecasting in SME Food and Beverage Retail}
\textbf{Author(s) \& Year:} Moch Azis Munawar, Ivan Diryana Sudirman (2026)\\
\textbf{Model/Method used:} LSTM, Naive Forecast, 7-hour Moving Average (MA7), Daily Seasonal Naive.\\
\textbf{Connection to our PS / gap:} Targets the forecast-to-action gap (3) using cost-based simulations.\\
\textbf{Literature Review Paragraph:}\\
The volatility of hourly sales in the food and beverage industry frequently causes inventory mismatches, leading to either costly spoilage or missed revenue. Munawar and Sudirman (2026) explored this issue by applying an LSTM network to forecast hourly ice cream sales in a culinary SME, benchmarking it against Naive, Moving Average, and Daily Seasonal Naive models. Their findings established that the LSTM model achieved the lowest MAE and RMSE by successfully capturing complex intra-day seasonal patterns that baseline models missed. Crucially, the authors extended their analysis to include cost-based simulations, demonstrating that LSTM-driven predictions actively reduce the financial penalties associated with overstocking and understocking. However, the authors explicitly acknowledged a major limitation: the model did not incorporate external drivers such as weather, holidays, or promotional events. LocalDemand directly addresses this recognized limitation. By utilizing Facebook Prophet configured with exogenous regressors for weather and holidays, LocalDemand captures the external demand shocks that purely autoregressive models miss, ultimately providing a more robust foundation for its automated restock recommendations.

\section{SME AI/ML Adoption}

\subsection{Optimizing Supply Chain Management for Small and Medium-Sized Businesses Using AI: A Comparative Analysis of GPT-4.5, Grok 3, Gemini 2.0, and DeepSeek R-1}
\textbf{Author(s) \& Year:} Shreeyash Vyakarnam et al. (2025)\\
\textbf{Model/Method used:} Generative AI (GPT-4.5, Grok 3, Gemini 2.0, DeepSeek R-1), XGBoost, Random Forest.\\
\textbf{Connection to our PS / gap:} Explores accessibility gap (2) using low-cost GenAI for SMBs.\\
\textbf{Literature Review Paragraph:}\\
As the capabilities of generative AI expand, its applicability to supply chain optimization for small and medium-sized businesses (SMBs) has garnered significant attention. Vyakarnam et al. (2025) conducted a comparative analysis of prominent generative models---including GPT-4.5, Grok 3, Gemini 2.0, and DeepSeek R-1---evaluating their capacity to execute demand planning and logistics tasks typically reserved for expensive ERP systems. While models like DeepSeek R-1 demonstrated promise in synthesizing demand planning insights from retail datasets, the authors found that all generative models failed to resolve complex, deterministic mathematical challenges such as the Vehicle Routing Problem. This reveals a critical limitation: while generative AI is highly accessible and cost-effective, it is fundamentally unsuited for rigorous numerical optimization in supply chains. This finding strongly supports the architectural decision behind LocalDemand. Rather than relying on generalized LLMs for mathematical forecasting, LocalDemand employs Facebook Prophet---a dedicated, purpose-built statistical engine---to ensure that inventory predictions are mathematically sound while maintaining the accessibility SMBs require.

\subsection{Technology Acceptance and Security Influences on AI Adoption among Small and Medium Enterprises}
\textbf{Author(s) \& Year:} Yuniarty et al. (2025)\\
\textbf{Model/Method used:} PLS-SEM on survey data based on TAM.\\
\textbf{Connection to our PS / gap:} Strongly reinforces accessibility gap (2) showing security and ease of use are critical over pure predictive power.\\
\textbf{Literature Review Paragraph:}\\
The technological superiority of an AI model does not guarantee its adoption by small business owners. Yuniarty et al. (2025) explored this phenomenon by applying the Technology Acceptance Model (TAM) to survey 131 SME owners in Indonesia regarding their intention to adopt AI tools. Utilizing Partial Least Squares–Structural Equation Modeling (PLS-SEM), the researchers discovered that perceived ease of use and perceived security were the primary drivers of technology adoption. Surprisingly, the perceived usefulness (the actual functional capability) of the AI did not significantly drive adoption intention if the system was deemed difficult to use or insecure. This finding is highly consequential for the design of forecasting tools in the SME sector, directly highlighting the technological accessibility gap. It confirms that providing a mathematically sound forecast is insufficient; the delivery mechanism must be friction-less. LocalDemand is architected around this exact behavioral insight. By abstracting the complex Facebook Prophet mathematics behind a simple, secure, and intuitive Flutter mobile interface, LocalDemand ensures that the barrier to entry for SME owners is minimized, driving actual adoption and practical utility.

\subsection{Technology Adoption of MSMEs in Food Supply Chain: The Role of AI in Traceability and Transparency}
\textbf{Author(s) \& Year:} Teerarat Aree, Pittawat Ueasangkomsate (2025)\\
\textbf{Model/Method used:} Linear regression analysis on survey data.\\
\textbf{Connection to our PS / gap:} Validates the accessibility gap (2) by confirming that MSMEs face severe financial and educational barriers to AI adoption.\\
\textbf{Literature Review Paragraph:}\\
The integration of AI into food supply chains offers significant operational benefits, yet its penetration into the micro, small, and medium enterprise (MSME) sector remains stalled by resource constraints. Aree and Ueasangkomsate (2025) investigated the role of AI in enhancing traceability and transparency among 92 MSMEs in the Thai food supply chain. Through linear regression analysis, the study confirmed that AI adoption significantly improves supply chain integrity and firm performance. Crucially, however, the authors identified a massive barrier to entry: MSMEs fundamentally lack both the financial capital and the technical "know-how" required to invest in and operate standard AI solutions. This directly validates the accessibility gap that LocalDemand seeks to bridge. The study underscores that enterprise-grade AI supply chain solutions are economically and operationally out of reach for micro-businesses. By delivering automated inventory forecasting through a low-cost, low-code mobile platform, LocalDemand democratizes access to AI, allowing SMEs to optimize their supply chains without requiring dedicated IT budgets or data science expertise.

\subsection{Factors influencing audit quality through the acceptance and use of AI by auditors in Thailand, affiliated with SMEs audit firms}
\textbf{Author(s) \& Year:} Kacha Somjaipeng, Watchareewan Jitsakul (2026)\\
\textbf{Model/Method used:} Technology Acceptance Model (TAM), PLS-SEM.\\
\textbf{Connection to our PS / gap:} Informs accessibility gap (2) from a behavioral standpoint.\\
\textbf{Literature Review Paragraph:}\\
Understanding the behavioral and environmental drivers behind SME technology adoption is critical for ensuring the successful deployment of AI solutions. Somjaipeng and Jitsakul (2026) utilized the Technology-Organization-Environment (TOE) framework and the Technology Acceptance Model (TAM) to evaluate AI adoption among 120 SME-affiliated auditors in Thailand. Using PLS-SEM analysis, they found that environmental factors---particularly regulatory frameworks---had a significant positive influence ($\beta = 0.562$) on the perceived benefits of AI, explaining over 70\% of the variance in AI acceptance. Although this study operates within the auditing sector rather than food service, it yields a crucial insight for the development of SME-focused technology: algorithmic superiority alone does not guarantee adoption. Solutions must be perceived as highly beneficial and environmentally aligned with the user's daily constraints. This principle underpins the design philosophy of LocalDemand, which addresses the SME accessibility gap by wrapping complex Prophet forecasting within an intuitive, highly accessible Flutter mobile interface that requires minimal technical expertise from the business owner.

\subsection{Predictive Marketing and Sales Analytics Using Machine Learning for Small and Medium E-Commerce Enterprises}
\textbf{Author(s) \& Year:} Swati Srivastava et al. (2026)\\
\textbf{Model/Method used:} Classification, regression, and ensemble learning framework proposition.\\
\textbf{Connection to our PS / gap:} Confirms accessibility gap (2) regarding budget and model complexity barriers.\\
\textbf{Literature Review Paragraph:}\\
The transition from reactive to predictive analytics remains a significant hurdle for smaller digital businesses. Srivastava et al. (2026) provided a comprehensive review of how machine learning---specifically regression, classification, and ensemble methods---can be leveraged by small and medium e-commerce enterprises for demand forecasting and targeted marketing. While the authors effectively argue that machine learning dramatically outperforms traditional statistical summaries by capturing dynamic consumer behavior, their primary contribution lies in identifying the barriers to entry. The study highlights that budget constraints, perceived model complexity, and inadequate data infrastructure prevent widespread ML adoption among SMEs. Although this research is largely theoretical and lacks empirical error metrics, it perfectly frames the "accessibility gap" that plagues the SME sector. It validates the core premise of LocalDemand: for advanced predictive analytics to be viable for small businesses, they must be abstracted away from complex data infrastructure and delivered through highly accessible, user-friendly mobile interfaces that require zero technical expertise to operate.

\section{Weather/External Factors}

\subsection{Enhanced sales forecasting of perishable and non-perishable goods using weather data with insights on Imports and price sensitivity}
\textbf{Author(s) \& Year:} Lalitha R et al. (2025)\\
\textbf{Model/Method used:} Random Forest and ARIMA.\\
\textbf{Connection to our PS / gap:} Partially addresses forecast-to-action gap (3) using predictions to drive dynamic pricing based on weather.\\
\textbf{Literature Review Paragraph:}\\
External environmental variables, particularly weather and market fluctuations, play a critical role in consumer dining behavior. Lalitha et al. (2025) integrated historical sales, local weather data, and ingredient costs into a machine learning framework to forecast restaurant demand for perishable and non-perishable goods. Utilizing Random Forest and ARIMA models, the authors achieved a high predictive accuracy (MSE = 53.97) and developed a dynamic pricing engine capable of adjusting menu prices in real-time based on fluctuating external drivers. While this approach successfully moves beyond mere prediction by implementing an actionable pricing strategy, it is highly reliant on continuous, real-time data streams that may be inaccessible to smaller enterprises. Furthermore, the actionability is focused entirely on revenue maximization through pricing rather than physical inventory control. LocalDemand pivots from this pricing-centric model to focus on the operational side, using similar external regressors (like weather) to generate accessible, automated restock recommendations for SMEs.

\subsection{Weather-Integrated Recommendations for Food Choices with SHAP Explanation}
\textbf{Author(s) \& Year:} Nikunj Jain et al. (2024)\\
\textbf{Model/Method used:} Random Forest, kNN, Decision Tree, Naïve Bayes, SVM, SHAP.\\
\textbf{Connection to our PS / gap:} Directly addresses forecast-to-action gap (3) and external regressors necessity.\\
\textbf{Literature Review Paragraph:}\\
Unpredictable weather fluctuations are a primary driver of demand volatility and subsequent food waste in the restaurant and cloud kitchen industry. Jain et al. (2024) addressed this challenge by developing a machine learning framework that correlates weather patterns with shifting consumer food preferences. Utilizing algorithms such as Random Forest and Support Vector Machines, augmented with SHAP for interpretability, the authors demonstrated that weather-integrated models could successfully generate adaptive menu recommendations. Their research highlighted a critical industry inefficiency: up to 40\% of predicted sales in the food sector fail to materialize due to unforeseen external factors like weather. This study strongly validates the necessity of incorporating exogenous variables into demand planning to close the forecast-to-action gap. However, while Jain et al. focused on recommending which dishes to serve, LocalDemand advances this concept into quantitative inventory management. By utilizing Facebook Prophet configured with specific weather regressors, LocalDemand moves beyond qualitative menu suggestions to generate exact, mathematically derived restocking quantities, directly translating weather-integrated demand predictions into actionable supply chain decisions for SME owners.

\subsection{Demand Forecasting in Supply Chain Management for Rossmann Stores using Weather Enhanced Deep Learning Model}
\textbf{Author(s) \& Year:} Nameer Ul Haq Qureshi et al. (2024)\\
\textbf{Model/Method used:} GRU deep learning model with Grid Search.\\
\textbf{Results:} GRU model outperformed LSTM, showing weather data improved accuracy.\\
\textbf{Literature Review Paragraph:}\\
Qureshi et al. (2024) presented a weather-enhanced deep learning model to improve demand forecasting in retail supply chain management. The study used historical sales data along with weather conditions, holidays, promotional activities, and store location information to build a forecasting model using the Gated Recurrent Unit (GRU) algorithm. The results showed that the proposed GRU model performed better than the LSTM model and produced more accurate sales forecasts by considering external factors such as weather. The study highlighted that better forecasting can help reduce overstocking and stock shortages while improving inventory planning. However, the proposed model was tested only on large retail stores and requires high computational resources, making it difficult to use directly in small businesses such as bakeries, cafés, and restaurants.

\subsection{AI-Based Sales Forecasting for F\&B Businesses}
\textbf{Author(s) \& Year:} Groene and Zakharov (date unspecified)\\
\textbf{Model/Method used:} XGBoost.\\
\textbf{Results:} Reduced RMSE by 22--33\% and MAE by 19--31\%.\\
\textbf{Literature Review Paragraph:}\\
Groene and Zakharov presented an AI-based sales forecasting approach to help small food and beverage (F\&B) businesses improve their operational and financial planning. The main objective of the study was to replace traditional experience-based forecasting with an intelligent machine learning model that can provide more accurate sales predictions. The authors developed an Extreme Gradient Boosting (XGBoost) model using more than five years of hourly sales data collected from a fast-food restaurant in Germany. The model was trained using historical sales data along with external factors such as weather conditions, holidays, local events, promotional activities, and macroeconomic indicators. The experimental results showed that the proposed AI model significantly outperformed manual forecasting methods by reducing the Root Mean Squared Error (RMSE) by 22\%--33\% and the Mean Absolute Error (MAE) by 19\%--31\%, leading to better inventory planning, reduced food waste, and improved operational efficiency. However, the study was conducted on a single fast-food franchise and required active involvement from managers during feature selection and model validation, which may limit its applicability to other businesses. Nevertheless, the research demonstrates the importance of integrating external factors into demand forecasting and strongly supports the development of AI-based inventory forecasting systems for small food businesses.

\subsection{Forecasting Restaurant Sales with Weather and Special Days using Facebook Prophet}
\textbf{Author(s) \& Year:} Güler et al. (date unspecified)\\
\textbf{Model/Method used:} Facebook Prophet.\\
\textbf{Results:} MAE = 0.76, MAPE = 0.18, RMSE = 0.82.\\
\textbf{Literature Review Paragraph:}\\
Güler et al. developed a restaurant sales forecasting system using the Facebook Prophet model to improve demand prediction. The study used around eight months of POS sales data along with weather information and holiday data to generate 30-day sales forecasts. The results showed that the model accurately captured seasonal trends and the effects of weather and holidays on sales, with the best performance achieving MAE = 0.76 and RMSE = 0.82. However, the research was based on data from only one restaurant and did not consider other factors such as promotions or local events. It also focused only on forecasting and did not provide a mobile application or inventory recommendations. This study is closely related to our project because it shows that using the Prophet model with weather and holiday data can improve demand forecasting. Our proposed system builds on this by offering a simple mobile app that provides daily demand forecasts and restocking recommendations for small bakeries, cafés, and restaurants.


\section{Inventory Optimization \& Restock}

\subsection{Evaluation of Optimized LSTM and Moving Average in SME Inventory Forecasting and Simulation}
\textbf{Author(s) \& Year:} Intan Rahmatillah et al. (2026)\\
\textbf{Model/Method used:} LSTM and Moving Average (MA7).\\
\textbf{Connection to our PS / gap:} Bridges forecast-to-action gap (3) using an inventory purchasing simulation.\\
\textbf{Literature Review Paragraph:}\\
While many forecasting studies stop at calculating error metrics, Rahmatillah, Herwanto, and Sudirman (2026) emphasize the necessity of translating predictive accuracy into operational inventory metrics. The authors compared an optimized LSTM network against a baseline 7-day moving average using daily sales data from an ice cream SME. The optimized LSTM not only achieved a superior predictive accuracy (RMSE = 1.65, MAPE = 45.06\%), but more importantly, it eliminated lost demand entirely when applied to a simulated weekly purchasing cycle. Despite generating a negligible overstock of 1.57 units, the study successfully validates the premise that deep learning can directly inform SME procurement strategies. However, the model relied on a rich dataset (nearly 1,000 days of history for a single item) and did not address the integration of external regressors. LocalDemand builds upon this operational simulation approach but adapts it for Prophet-based predictions, integrating weather data and specifically designing around the cold-start limitations where such extensive historical data is unavailable.

\subsection{Demand Forecasting System, Small Business Optimization}
\textbf{Model/Method used:} SARIMA, DeepAR, and Hybrid.\\
\textbf{Results:} DeepAR achieved highest forecasting accuracy of 81.0\%.\\
\textbf{Literature Review Paragraph:}\\
This paper presents a demand forecasting system designed for small businesses using SARIMA, DeepAR, and a hybrid forecasting approach to improve inventory management. The authors compare traditional statistical forecasting with deep learning and report that DeepAR achieved the highest forecasting accuracy of 81.0\%, while SARIMA achieved 50.76\% and the hybrid model 77.14\%. The proposed system also provides a web interface for uploading sales data and generating demand forecasts, enabling better stock planning and reducing overstocking and stock shortages. However, the model does not incorporate external factors such as weather, holidays, or promotions. This work is closely related to our proposed system because it demonstrates the importance of affordable demand forecasting tools for small businesses, while our work further enhances this concept by using the Prophet forecasting model with external factors to generate simple restock recommendations through a mobile application.

\subsection{AI Predictive Analytics, Sustainable Restaurant Ops (AID-PAF)}
\textbf{Model/Method used:} Temporal-Fusion Neural Architecture (TFNA) + Waste-Aware LP.\\
\textbf{Results:} Achieved MAPE of 6.5\%, 40\% reduction in food waste.\\
\textbf{Literature Review Paragraph:}\\
This paper proposes an AI-Driven Predictive Analytics Framework (AID-PAF) for sustainable restaurant operations by integrating a Temporal-Fusion Neural Architecture (TFNA) with a waste-aware linear programming model for inventory optimization. The forecasting model incorporates historical sales along with external factors such as weather, holidays, promotions, and social media sentiment to generate accurate demand predictions. Experimental evaluation using an AnyLogic-MATLAB digital twin simulator achieved a MAPE of 6.5\%, 96.2\% customer fill rate, 40\% reduction in food waste, and 21\% reduction in procurement costs, demonstrating significant improvements in operational efficiency and sustainability. However, the framework relies on large datasets and computationally intensive deep learning models, making it less suitable for small businesses with limited resources. This work is closely related to our proposed system because it highlights the importance of combining demand forecasting with inventory optimization to reduce food waste, while our solution focuses on providing an affordable mobile application using the Prophet forecasting model with weather and holiday data for small food businesses.

\subsection{Smart Stock: Small Warehouse Inventory}
\textbf{Model/Method used:} XGBoost-based demand forecasting.\\
\textbf{Results:} 96\% classification accuracy and R$^2$ > 0.95.\\
\textbf{Literature Review Paragraph:}\\
This paper proposes a Smart Inventory Management System for small and medium-sized enterprises (SMEs) that integrates XGBoost-based demand forecasting, QR code product identification, Raspberry Pi edge devices, IoT, and cloud storage to automate inventory monitoring and replenishment decisions. The proposed system achieved 96\% classification accuracy and an R$^2$ value greater than 0.95, outperforming traditional machine learning models such as Random Forest, Linear Regression, and KNN. Real-time stock monitoring and automated alerts help reduce stockouts, overstocking, and manual errors while providing a cost-effective alternative to expensive ERP systems such as SAP and Oracle. However, the system is primarily designed for warehouse inventory management and does not incorporate external demand factors such as weather or holidays. This work is relevant to our proposed system because it highlights the benefits of AI-powered inventory management for SMEs, while our project extends this concept by focusing on small food businesses, using the Prophet forecasting model and external factors to generate accurate daily inventory and restocking recommendations.

\subsection{Predictive Analytics, Food Supply Chain \& Waste}
\textbf{Model/Method used:} LSTM, XGBoost, and Random Forest.\\
\textbf{Results:} XGBoost achieved 94.3\% accuracy, 53.8\% food waste reduction.\\
\textbf{Literature Review Paragraph:}\\
This paper investigates the application of machine learning models including LSTM, XGBoost, and Random Forest for demand forecasting, inventory optimization, and food waste reduction in restaurant supply chains. Using restaurant POS data, inventory logs, and waste records collected over three months, the authors compared the performance of the three models and found that XGBoost achieved the best results with 94.3\% accuracy, 95.1\% fill rate, 53.8\% food waste reduction, and 27.6\% cost savings, outperforming LSTM and Random Forest. The study demonstrates that predictive analytics can significantly improve inventory planning, reduce waste, and enhance restaurant profitability. However, the research is limited to a comparison of machine learning models and does not evaluate Prophet or provide a lightweight mobile solution for small businesses. This work is highly relevant to our proposed system because it validates the importance of AI-driven demand forecasting for inventory optimization, while our project focuses on an affordable Prophet-based mobile application that incorporates weather and holiday information to assist small food businesses in making daily restocking decisions.


\section{Mobile/App-Based Decision Support}

\subsection{Mobile application for the intelligent management of small shops using the GPT-4 model}
\textbf{Author(s) \& Year:} Gonzalo Sakuda et al. (2024)\\
\textbf{Model/Method used:} GPT-4 LLM.\\
\textbf{Connection to our PS / gap:} Addresses accessibility gap (2) perfectly via a mobile-first AI tool, though it lacks robust empirical time-series forecasting.\\
\textbf{Literature Review Paragraph:}\\
The integration of Artificial Intelligence into mobile platforms offers a promising avenue to enhance the operational capabilities of small retailers. Sakuda, Leon, and Rojas (2024) developed an Android application integrating the GPT-4 Large Language Model to act as an intelligent management assistant for small shops in Peru. Their system allowed users to query real-time sales data and receive conversational recommendations on inventory and promotions. A qualitative assessment revealed strong user approval, with 66.67\% of surveyed shop owners rating the tool highly for usability. While this effectively demonstrates how to bridge the technological accessibility gap for small businesses, the study relies exclusively on generative AI for its insights rather than empirical forecasting models. It lacks a rigorous mathematical foundation for inventory predictions, relying instead on contextual LLM outputs. LocalDemand addresses this limitation by marrying the accessible mobile-app paradigm demonstrated here with the mathematical robustness of Facebook Prophet, ensuring that recommendations are driven by precise time-series calculations rather than generalized generative AI.

\subsection{A Mobile Application for Inventory Management to Support Digital Transformation in Small Retail Businesses}
\textbf{Author(s) \& Year:} Thanakrit Janchidfah et al. (2026)\\
\textbf{Model/Method used:} Technology Acceptance Model (TAM).\\
\textbf{Connection to our PS / gap:} Addresses accessibility gap (2) by demonstrating small retailers need lightweight mobile tools over complex ERPs.\\
\textbf{Literature Review Paragraph:}\\
In addressing the digital transformation needs of small and medium enterprises (SMEs), Janchidfah et al. (2026) demonstrate the efficacy of lightweight, user-centered mobile applications over complex enterprise resource planning systems. The authors developed "Stock Easy," an inventory management tool, and evaluated it using the Technology Acceptance Model (TAM) with 132 retail stakeholders. Their findings indicate that perceived ease of use (mean = 4.52) and process clarity are critical drivers of user satisfaction, which in turn strongly correlates with an organization's readiness for digital transformation (r = 0.76). While the study successfully highlights the necessity of accessible, mobile-first solutions for small retailers, the application remains functionally reactive. It relies on manual stock tracking and lacks predictive analytics. This highlights a significant opportunity for our proposed application, LocalDemand, which will build upon this accessible mobile framework by integrating proactive, machine-learning-driven demand forecasting.

\section{Base Paper Recommendation}

\textbf{Primary Methodological Base Paper:}\\
\textit{Forecasting Restaurant Sales with Weather and Special Days using Facebook Prophet} (Güler et al.) OR \textit{State-Level Drug Retail Sales Forecasting Using XGBoost and Facebook Prophet} (Kamboj et al., 2025). These share our core model (Prophet) combined with weather/holiday regressors.\\

\textbf{Primary Application Context Base Paper:}\\
\textit{Comparative Evaluation of SARIMA, FB Prophet, and LSTM for Hourly Demand Forecasting in Saudagar Kopi} (Aqila et al., 2026) strongly matches the small food business (SME) environment.

\end{document}
